\begin{figure}[tbp]
    \centering

    \begin{tikzpicture}[
        >=Latex,
        font=\small,
        update/.style={
            -{Latex[length=2.5mm]},
            very thick
        }
    ]

    % ================================================================
    % TOP: geometric schematic
    % ================================================================

    \coordinate (x) at (0,0);
    \def\R{1.7}

    % L2 ball
    \draw[thick] (x) circle (\R);

    % Starting point
    \fill (x) circle (2pt);
    \node[below left=2pt] at (x) {$x$};


    % ----------------------------------------------------------------
    % Radius epsilon
    % ----------------------------------------------------------------
    \draw[
        {Latex[length=1.7mm]}-{Latex[length=1.7mm]}
    ]
        (x) -- ++(0,-\R)
        node[midway, left=4pt] {$\epsilon$};

    \node[below=5pt] at (0,-\R)
        {$L_2$ ball};


    % ----------------------------------------------------------------
    % Normalised update
    %
    % Deliberately exaggerated so that it remains visible.
    % ----------------------------------------------------------------
    \draw[update]
        (x) -- ++(0.85,0.85)
        coordinate (normalend);

    % Put the label ABOVE the endpoint rather than on the arrow.
    \node[
        above=5pt of normalend,
        align=center
    ] {
        \textbf{normalised}\\[-1pt]
        $\|u\|_2=0.1\epsilon$
    };


    % ----------------------------------------------------------------
    % Released update
    % ----------------------------------------------------------------
    \draw[update]
        (x) -- ++(4.5,0)
        coordinate (releaseend);

    % Put this label well above the horizontal arrow.
    \node[
        above=6pt of releaseend,
        anchor=south east,
        align=center
    ] {
        \textbf{released}\\[-1pt]
        $\|u\|_2=6.4\epsilon$
    };

    % Not-to-scale note BELOW the long arrow.
    \node[
        below=6pt of releaseend,
        anchor=north east,
        font=\scriptsize,
        text=gray
    ] {
        arrow lengths not to scale
    };


    % ================================================================
    % BOTTOM: algebra
    %
    % Crucially, this is BELOW the entire geometric diagram.
    % Nothing is competing for space around the circle.
    % ================================================================

    \coordinate (eqtop) at (2.25,-2.65);


    % ----------------------------------------------------------------
    % Normalised equation
    % ----------------------------------------------------------------
    \node[
        anchor=north east,
        align=center,
        text width=5.2cm
    ] (normalbox) at ($(eqtop)+(-0.35,0)$) {
        \textbf{Dimension-normalised}\\[5pt]
        $\displaystyle
        \eta=\frac{\lambda\epsilon}{\sqrt d}
        $\\[5pt]
        $\displaystyle
        \|u\|_2
        =\eta\sqrt d
        =\lambda\epsilon
        =0.1\epsilon
        $
    };


    % ----------------------------------------------------------------
    % Released equation
    % ----------------------------------------------------------------
    \node[
        anchor=north west,
        align=center,
        text width=5.2cm
    ] (releasebox) at ($(eqtop)+(0.35,0)$) {
        \textbf{Released}\\[5pt]
        $\displaystyle
        \eta=\lambda\epsilon
        $\\[5pt]
        $\displaystyle
        \|u\|_2
        =\eta\sqrt d
        =\lambda\epsilon\sqrt d
        =6.4\epsilon
        $
    };


    % ----------------------------------------------------------------
    % 64x punchline
    %
    % Positioned BELOW both derivations, so it cannot overlap either.
    % ----------------------------------------------------------------
    \node[
        draw,
        rounded corners=2pt,
        thick,
        inner xsep=12pt,
        inner ysep=6pt,
        below=9mm of normalbox.south east,
        anchor=north,
        align=center
    ] {
        $\sqrt d=\sqrt{4096}=64$
        \qquad $\Longrightarrow$ \qquad
        \textbf{$64\times$ larger update}
    };

    \end{tikzpicture}

    \caption{
        Effect of step-size normalisation on the unprojected update.
        With $\lambda=0.1$ and $d=4096$, dimension normalisation gives
        an update norm of $\lambda\epsilon=0.1\epsilon$.
        The released step size omits the $1/\sqrt d$ factor, increasing
        the unprojected update norm by $\sqrt d=64$ to
        $\lambda\epsilon\sqrt d=6.4\epsilon$.
        Arrow lengths are not to scale.
    }
    \label{fig:step-size-schematic}
\end{figure}